\documentclass[lettersize,journal]{IEEEtran}
\usepackage{amsmath,amsfonts}
\usepackage{algorithm}
% \usepackage{algorithmic}
\usepackage{algpseudocode}
\renewcommand{\Comment}[1]{\hfill$\triangleright$ #1}
\usepackage{array}
%\usepackage[caption=false,font=normalsize,labelfont=sf,textfont=sf]{subfig}
\usepackage[]{subcaption}
\usepackage{arydshln}
\usepackage{textcomp}
\usepackage{stfloats}
\usepackage{url}
\usepackage{verbatim}
\usepackage{graphicx}
\usepackage{IEEEtrantools}
\usepackage{hyperref}
\usepackage{setspace}
\usepackage{soul}
\usepackage{xcolor}
\usepackage{makecell} % added for compact multi line table headers, Related Work comparison table
\usepackage{amsthm} % added for the formal lemma, corollary, and remark environments in the appendix
\theoremstyle{plain}
\newtheorem{lemma}{Lemma}
\newtheorem{corollary}[lemma]{Corollary}
\theoremstyle{remark}
\newtheorem{remark}{Remark}
\hyphenation{op-tical net-works semi-conduc-tor IEEE-Xplore}
\def\BibTeX{{\rm B\kern-.05em{\sc i\kern-.025em b}\kern-.08em
		T\kern-.1667em\lower.7ex\hbox{E}\kern-.125emX}}
\usepackage{balance}
\setlength{\textfloatsep}{8pt}
\newcommand{\orcidlink}[1]{\href{https://orcid.org/#1}{\includegraphics[scale=0.06]{pictures/orcid.png}}}
\newcommand{\TBD}{\textbf{TBD}} % placeholder for unrun quantitative results, replace once measured

\begin{document}
	\title{NEXUS: Unified Post-Quantum Zero-Trust Detection of Selective Time Delay and Hidden Forwarding Attacks via Federated LSTM and Controller-Failover Blockchain in SDVN}
	\author{
		% TODO: author block pending confirmation of the final author list for this
		% paper. Placeholder below follows the template structure only.
    \IEEEauthorblockN{Patikiri Arachchige Don Shehan Nilmantha Wijesekara \textsuperscript{\orcidlink{0000-0002-3045-1596}}\\
    Pathiraja Mudiyanselage Sajani Chamindi Bandara Gunarathna 
    \textsuperscript{\orcidlink{0009-0000-8230-9631}}\\
    Lankahaluge Nipuni Kaushalya Fernando 
    \textsuperscript{\orcidlink{0009-0003-3888-0015}} \\
    Siyagunakosgodage Apsari Udithara Fernando \textsuperscript{\orcidlink{0009-0008-1328-8855}} }
		
		\thanks{P.A.D.S.N. Wijesekara (nilmantha@eie.ruh.ac.lk), are with Dept. Electr. \& Informa. Eng., Faculty0009-0003-3888-0015 of Engineering, University of Ruhuna, Galle 80000, Sri Lanka.}
	}
	\markboth{IEEE Transactions on Intelligent Transportation Systems,~Vol.~XX, No.~XX, 2027}
	{P.A.D.S.N. Wijesekara \MakeLowercase{\textit{et al.}}}
	
	%\IEEEpubid{0000--0000/00\$00.00~\copyright~2027 IEEE.}
	
	\maketitle
	
	\begin{abstract}
		The centralized programmability of software defined vehicular networks
		exposes safety critical traffic to selective time delay and hidden
		forwarding attacks, which desynchronize collision avoidance by delaying
		priority packets or covertly duplicate sensitive data to unauthorized
		destinations while mimicking legitimate traffic to evade volume based
		detection. Existing defenses cover only one attack class, assume the
		controller itself is trustworthy, or centralize raw traffic that
		exposes driver location; none detects both jointly under zero trust
		extended to the controller. This paper presents NEXUS, spanning a dual
		mode detection engine, a post quantum cryptographic layer of lattice
		based signatures and zero knowledge proofs, a federated classifier
		across roadside units that never centralizes traffic, and a witness
		mechanism for passive attacks that leave no cryptographic or
		volumetric trace. A flat, registered controller set with a delay
		evidence trust channel closes a gap in endorsement only scoring,
		penalizing a timing manipulating controller even when every flow rule
		it issues is authorized. Evaluated with all eight variants active
		simultaneously on a Manhattan trace, NEXUS sustains macro Matthews
		correlation coefficient \TBD\ against \TBD\ for the strongest baseline
		at \TBD\ percent penetration, while every baseline shows the coverage
		collapse this joint setting is designed to expose; end to end latency
		stays within the hundred millisecond bound to four hundred vehicles.
		These results show that joint detection across both attack classes,
		extended to the controller itself, is achievable without the false
		positive or latency cost a naive combination of single class defenses
		would incur.
	\end{abstract}
	
	\begin{IEEEkeywords}
		Software Defined Vehicular Networks, Post Quantum Cryptography, Methods for security and privacy, Artificial Intelligence, Vehicular Ad hoc Networks, Connected and Autonomous Vehicles.
	\end{IEEEkeywords}
	
	
	\section{Introduction}\label{introduction}
	
	\subsection{Background}\label{subsec:background}
	
	\IEEEPARstart{S}{oftware} defined vehicular networks decouple traffic
	forwarding at roadside units from routing decisions made by a logically
	centralized controller, enabling flexible flow management across a
	rapidly changing vehicular topology. This same programmability creates
	an attack surface absent from static networks: a compromised controller,
	roadside unit, or vehicle can selectively delay safety critical packets
	such as collision avoidance warnings to desynchronize dependent systems,
	or can duplicate and reroute sensitive data to an unauthorized
	destination while the original packet still reaches its legitimate
	recipient unmodified, evading any check based on whether delivery
	occurred at all. Both attack classes are engineered to mimic legitimate
	traffic closely enough that volume based and static threshold detection
	fail to distinguish them from ordinary behaviour, and vehicular
	mobility, through handoff induced latency jitter and rapid topology
	change, widens this camouflage beyond what any static deployment would
	need to tolerate.
	
	\subsection{Limitations of Existing Work}\label{subsec:intro_limitations}
	
	Prior defenses close only part of this problem. Frameworks such as SIFT
	and FSDM assume the controller issuing flow rules is itself trustworthy,
	so a controller that becomes the adversary defeats them by construction.
	Path validation approaches such as HSA centralize raw packet metadata at
	a collector, exposing exactly the location and movement data a
	vehicular deployment must protect. Timing based protocols such as TAP
	rely on a single static threshold that a slow, deliberately paced
	attacker can evade, and every one of these lines of work addresses only
	selective time delay or only hidden forwarding, never both under one
	architecture. None combines controller inclusive zero trust, post
	quantum attestation, and federated detection, the gap
	Table~\ref{tab:related_work} of Section~\ref{related} sets out in full.
	
	\subsection{Problem Statement}\label{subsec:problem_statement}
	
	The practical consequence of this gap is concrete rather than abstract.
	A vehicle that brakes late because its collision warning was
	selectively delayed past the safety window, and a vehicle whose route
	and destination history is silently duplicated to an unauthorized
	listener through a covertly rerouted flow, are physical and privacy
	harms that can arise from the same compromised entity at the same time.
	A defense that watches for only one attack class, or that extends trust
	to the controller issuing its own flow rules, leaves the other harm to
	proceed entirely undetected while the first is being handled, and does
	so precisely under the high mobility conditions where both attacks are
	hardest to see.
	
	\subsection{Research Questions}\label{subsec:research_questions}
	
	This paper is organised around three research questions that the
	experiments of Section~\ref{results} answer directly.
	\begin{itemize}
		\item \textbf{RQ1:} Does a framework built for joint detection sustain
		macro averaged detection quality as simultaneous, multi class attack
		penetration increases, where single class defenses evaluated under
		the same joint traffic are expected to fail even on the class each
		was designed for?
		\item \textbf{RQ2:} Do the post quantum cryptographic layer and the
		blockchain consensus pipeline it depends on sustain end to end
		latency within the safety critical bound as network scale grows,
		given that every packet now carries a lattice based signature and
		every flow modification requires roadside unit endorsement?
		\item \textbf{RQ3:} Is each architectural layer, the cryptographic
		layer, the federated classifier, and the controller inclusive zero
		trust architecture, load bearing, such that removing any one of them
		produces a proportionate loss of the security guarantee it was added
		to provide?
	\end{itemize}
	
	\subsection{Contributions}\label{subsec:contributions}
	
	The principal contributions of this paper are:
	\begin{itemize}
		\item \textbf{Mobility aware joint attack modeling:} the formal
		definition of eight attack variants spanning both selective time
		delay and hidden forwarding, characterising how vehicular mobility
		amplifies each one beyond what static deployments experience.
		\item \textbf{Dual mode detection engine:} a lightweight, rule based
		detector deployable on resource constrained on board units, paired
		with a full pipeline at roadside units that adds cryptographic,
		witness, and classifier evidence without either mode sacrificing
		fidelity at the other's expense.
		\item \textbf{Unified post quantum cryptographic layer:} per hop
		lattice based signatures with randomized batch verification,
		combined with zero knowledge proofs for timing compliance and next
		hop legitimacy, designed specifically for vehicular data plane and
		flow rule integrity.
		\item \textbf{Federated classifier for passive detection:} a
		recurrent model trained across roadside units without centralizing
		raw traffic, whose classification head is trained on injection side
		ground truth independent of the features it consumes.
		\item \textbf{Multi controller zero trust architecture:} a flat,
		registered controller set with blockchain enforced trust scoring and
		a delay evidence channel that penalizes a timing manipulating
		controller even when every flow rule it issues remains individually
		authorized.
		\item \textbf{Witness based passive attack detection:} a broadcast
		radio monitoring mechanism that reaches a Byzantine fault tolerant
		quorum of independently verified alerts before any trust penalty is
		applied, covering attacks that leave no cryptographic or volumetric
		trace.
	\end{itemize}
	
	\subsection{Novelty}\label{subsec:novelty}
	
	This work makes four individually novel contributions that are further
	amplified in combination. The formal modeling of eight selective time
	delay and hidden forwarding attack variants, explicitly characterising
	how vehicular mobility amplifies each one, is absent from prior
	literature, which addresses these attacks only in static networks or
	treats the two classes separately. The dual mode detection architecture,
	pairing lightweight rule based signatures with federated classification
	and broadcast radio witness monitoring within one adaptive framework,
	is the first to cover the full spectrum of vehicular node capability
	without sacrificing fidelity against passive covert variants that
	produce no observable network anomaly. The integration of per packet
	lattice based signatures, randomized batch verification, and
	transparent zero knowledge proofs for both timing compliance and next
	hop legitimacy into one post quantum cryptographic layer is novel in
	both scope and target domain. The flat, multi controller architecture,
	with blockchain enforced trust scoring, distributed key generation
	among roadside unit ring members, and smart contract triggered
	failover, is the first controller independent design in which the
	controller is excluded from the key plane entirely. Collectively, this
	is the first framework to jointly detect and mitigate both attack
	classes under a zero trust, controller independent architecture that
	simultaneously preserves driver privacy and satisfies vehicular latency
	constraints.
	
	The remainder of this paper is organised as follows. Section~\ref{related}
	reviews related work and positions the gap this paper addresses.
	Section~\ref{system_model} formalises the network model, threat model,
	and eight attack variants. Section~\ref{framework} presents the NEXUS
	framework. Sections~\ref{evaluation_setup} and~\ref{results} describe
	the evaluation setup and report the three experiments. Section~\ref{discussion}
	discusses the results, and Section~\ref{conclusion} concludes.
	
	
	\section{Related Work}\label{related}
	
	Selective time delay and hidden forwarding attacks in software defined
	vehicular networks have largely been studied as separate problems, each
	addressed by a body of work built around its own threat model and its own
	notion of what the controller is permitted to do. Prior defenses for
	selective time delay concentrate on control plane saturation and slow
	table exhaustion, prior defenses for hidden forwarding concentrate on path
	validation and flow conservation, and the two lines of work rarely share
	an evaluation setting, let alone a common architecture. Table~\ref{tab:related_work}
	places eight representative directions against six properties that
	together determine whether a defense is suitable for the setting
	considered in this study: which attack class the work targets, what trust
	it places in the controller, whether it uses post quantum cryptography,
	whether detection is federated across edge nodes rather than centralized,
	whether the design accounts for vehicular mobility, and whether the two
	attack classes are detected jointly rather than in isolation.
	
	\begin{table*}[!t]
		\centering
		\caption{Comparison of representative defenses against selective time delay and hidden forwarding attacks in software defined vehicular networks.}
		\label{tab:related_work}
		\renewcommand{\arraystretch}{1.25}
		\footnotesize
		\begin{tabular}{|>{\raggedright\arraybackslash}p{2.0cm}|>{\raggedright\arraybackslash}p{3.1cm}|>{\raggedright\arraybackslash}p{2.6cm}|>{\centering\arraybackslash}p{1.0cm}|>{\centering\arraybackslash}p{1.3cm}|>{\centering\arraybackslash}p{1.6cm}|>{\centering\arraybackslash}p{1.7cm}|}
			\hline
			\textbf{Work} & \textbf{Attack Class Covered} & \textbf{Controller Trust} & \makecell{\textbf{Post}\\\textbf{Quantum}} & \textbf{Fed.} & \makecell{\textbf{Mobility}\\\textbf{Aware}} & \makecell{\textbf{Joint}\\\textbf{TD+HF}} \\
			\hline
			SIFT~\cite{Pascoal2017} & Time delay (TCAM exhaustion) & Assumed trusted & No & No & No & No \\
			\hline
			TAP~\cite{Arsalan2018} & Time delay (selective packet delay) & Assumed trusted & No & No & Yes & No \\
			\hline
			HSA~\cite{Mohammadi2016} & Hidden forwarding (path deviation) & Assumed trusted & No & No & No & No \\
			\hline
			FADE~\cite{Zhang2021FADE} & Hidden forwarding (duplication) & Assumed trusted & No & No & No & No \\
			\hline
			TFMD SDVN~\cite{Nayak2022} & Generic misbehavior & Semi trusted & No & No & Yes & No \\
			\hline
			Liu et al.~\cite{Liu2021BFLVEC} & Generic intrusion & Assumed trusted & No & Yes & Yes & No \\
			\hline
			FTISCON~\cite{Krishnan2023} & Flow rule tampering & Blockchain verified\textsuperscript{*} & No & No & No & No \\
			\hline
			Post quantum V2R authentication~\cite{ElDalahmeh2025} & Not attack detection (key exchange) & Assumed trusted & Yes & No & Yes & No \\
			\hline
			\textbf{Proposed (NEXUS)} & \textbf{Time delay and hidden forwarding, jointly} & \textbf{Zero trust, controller included} & \textbf{Yes} & \textbf{Yes} & \textbf{Yes} & \textbf{Yes} \\
			\hline
		\end{tabular}
		\vspace{2pt}
		\begin{flushleft}
			\footnotesize
			\textsuperscript{*}Verifies that an installed rule matches a committed policy; does not itself withdraw trust from the controller that issued the rule.
		\end{flushleft}
	\end{table*}
	
	No entry in Table~\ref{tab:related_work} combines the last three columns.
	TFMD SDVN and the framework of Liu et al. are mobility aware and, in the
	latter case, train collaboratively across distributed edge nodes, yet
	neither targets selective time delay or hidden forwarding specifically,
	and neither questions the honesty of the controller itself. Post quantum
	vehicular authentication secures the exchange between a vehicle and the
	infrastructure but performs no attack detection at all, and the
	blockchain based flow rule verification of FTISCON checks a rule against
	policy without extending distrust to the entity that issued the rule in
	the first place. The joint column is empty for every prior work, which is
	the gap this study closes: a single framework that treats the controller
	as untrusted along with every vehicle and roadside unit, detects both
	attack classes at once rather than in isolation, and does so using post
	quantum attestation and privacy preserving federated learning rather than
	a centrally trained model over raw traffic.
	
	
	\section{System Model and Attack Formalization}\label{system_model}
	
	\subsection{Network and Threat Model}\label{subsec:network_threat}
	
	The network comprises a mobile vehicle population, a fixed grid of
	roadside units, and a flat, registered set of controllers, connected as
	formalised in Sections~\ref{subsec:ctrl_trust} and~\ref{subsec:blockchain}.
	Controllers issue flow modification instructions that roadside units
	install and enforce, and traffic moves over both vehicle to
	infrastructure and vehicle to vehicle paths.
	
	The threat model assumes an adversary that can delay, duplicate,
	fabricate, or covertly forward packets, and that can exploit flow
	modification installation and forwarding behaviour from either the
	control plane or the data plane, never both at once from a single
	compromised entity. Zero trust extends to the controller itself: every
	controller is granted read only blockchain access, and its behaviour is
	audited only through independent roadside unit logging rather than any
	self reporting. Roadside units are individually corruptible but
	collectively honest under a Byzantine fault tolerance bound of
	$f<n/3$, and no controller in the registered set is permanently active
	or inherently trusted. This one bound $f$ is the parameter behind every
	threshold expressed as $f{+}1$ or $2f{+}1$ later in this paper, whether
	the quorum being reached is roadside unit endorsement, controller trust
	evidence, or witness confirmation, rather than a separate value chosen
	independently for each mechanism.
	
	\subsection{Attack Scenarios}\label{subsec:attack_scenarios}
	
	Two attack classes are addressed, each exploiting the same centralized
	logic from a different vantage point. Selective time delay identifies
	safety critical packets and introduces an intentional, variable lag
	before forwarding, so packets arrive late rather than being dropped and
	dependent systems act on stale data. Figure~\ref{fig:std_attacks} shows
	the four variants: control plane and data plane delay injection, and
	control plane and data plane slow flow table exhaustion, the latter
	pair achieving the same delay by starving the roadside unit's
	forwarding table rather than by injecting delay directly.
	\begin{figure}[!t]
		\centering
		\includegraphics[width=0.95\linewidth]{selective_time_delay.png}
		\caption{Selective Time Delay Attack Variants}
		\label{fig:std_attacks}
	\end{figure}
	Hidden forwarding instead leaves the original packet on its legitimate
	path while a duplicate is secretly routed to an unauthorized listener,
	enabling eavesdropping without any observable disruption to normal
	delivery. Figure~\ref{fig:hf_attacks} shows the four variants, split by
	adversary plane as above and by whether the duplicate is actively
	fabricated with a forged signature or passively retransmitted with the
	original, genuinely valid one.
	\begin{figure}[!t]
		\centering
		\includegraphics[width=0.95\linewidth]{Hidden Forwarding Attack Versions1.png}
		\caption{Hidden Forwarding Attack Variants}
		\label{fig:hf_attacks}
	\end{figure}
	Table~\ref{tab:taxonomy} summarises all eight variants against the
	plane each occupies, the property that makes each one covert, and the
	primary detector Section~\ref{framework} assigns to it.
	\begin{table}[!t]
		\centering
		\caption{Attack Taxonomy}
		\label{tab:taxonomy}
		\renewcommand{\arraystretch}{1.2}
		\footnotesize
		\begin{tabular}{|c|c|>{\raggedright\arraybackslash}p{2.7cm}|>{\raggedright\arraybackslash}p{2.3cm}|}
			\hline
			\textbf{Variant} & \textbf{Plane} & \textbf{Covertness} & \textbf{Primary Detector} \\
			\hline
			S1 & Control & Delay mimics handoff jitter & Signature S1 \\
			\hline
			S2 & Data & Delay mimics handoff jitter & Signature S2 \\
			\hline
			S3 & Control & Installation rate within legitimate distribution & Signature S3 \\
			\hline
			S4 & Data & Installation rate within legitimate distribution & Signature S4 \\
			\hline
			S5 & Control & Original path unaffected, duplicate fabricated & Signature S5 \\
			\hline
			S6 & Data & Original path unaffected, duplicate fabricated & Signature S6 \\
			\hline
			S7 & Control & Duplicate is genuinely authentic, passes all checks & Classifier, witness quorum \\
			\hline
			S8 & Data & Duplicate is genuinely authentic, passes all checks & Classifier, witness quorum \\
			\hline
		\end{tabular}
	\end{table}
	
	\subsection{Mobility Induced Threat Amplification}\label{subsec:mobility_amplification}
	
	Vehicular mobility is not incidental to these eight variants, it
	structurally widens the camouflage available to each one. Selective
	time delay is masked because legitimate handoff latency during flow
	rule recomputation falls in the same range an attacker would inject,
	so a detector calibrated as though $\delta_h{=}0$ is blind to exactly
	the overlap realistic mobility produces. Hidden forwarding is masked
	for a different reason: flow conservation detectors that count packets
	over a flow's lifetime need enough samples to test reliably, and a
	vehicle's residence time inside a single roadside unit's zone bounds
	how many samples such a detector ever gets to see. This condition is
	given in Equation~\eqref{eq:observation_window}.
	\begin{equation}
		N_{obs} = r \cdot t_{zone} < N_{min}(\alpha)
		\label{eq:observation_window}
	\end{equation}
	As given in Equation~\eqref{eq:observation_window}, $N_{obs}$ is the
	number of packets observable during a single zone crossing, the
	product of the per flow packet rate $r$ and the zone residence time
	$t_{zone}$, typically $9$ to $22$s at highway speed. Once this falls
	below the minimum sample size $N_{min}(\alpha)$ a flow conservation
	test needs at significance level $\alpha$, detectors such as FADE and
	HSA of Section~\ref{related} become statistically unreliable
	regardless of how the threshold is tuned. This is what motivates
	per packet, per hop cryptographic verification over flow level
	statistics as the primary hidden forwarding mechanism in
	Section~\ref{subsec:crypto_layer}, since a zero knowledge proof
	requires no observation window at all, complemented by the federated
	classifier of Section~\ref{subsec:fed_lstm} for the residual passive
	case a proof alone cannot cover.
	
	
	\section{NEXUS Framework}\label{framework}
	
	\subsection{Zero Trust Multi Controller Architecture}\label{subsec:ctrl_trust}
	
	NEXUS replaces the single trusted controller assumed by prior defenses
	with a flat, registered controller set in which no controller is
	inherently trusted and every controller is granted only read only access
	to the blockchain. Roadside units are dynamically assigned to whichever
	registered controller currently holds sufficient on chain trust, and a
	controller whose trust falls below the minimum threshold is revoked and
	its roadside units reassigned, with no pre designated backup that could
	itself be compromised.
	
	A structural gap in endorsement based trust scoring motivates the central
	architectural contribution of this layer. Conflict evidence is the count
	of roadside units that report a flow modification absent from the
	blockchain committed policy of Section~\ref{subsec:blockchain}, the same
	unauthorized installation Signature S3 of
	Section~\ref{subsec:dual_mode} tests for; a controller that manipulates
	packet timing while issuing topologically valid flow modification
	messages accumulates zero conflict evidence by this definition, since
	nothing about the installed rule itself is unauthorized, and would
	therefore be rewarded rather than penalized under an endorsement check
	alone. NEXUS closes this gap with a second, independent evidence channel
	keyed to Signature S1. The delay evidence count for controller $c_i$ at
	time $t$ is given in Equation~\eqref{eq:delay_evidence}.
	\begin{equation}
		\mathcal{E}_{delay}(c_i,t) =
		\left|\left\{
		r_k \in \mathcal{R}_{c_i}(t) :
		f_{S1}(r_k,t) = 1
		\right\}\right|
		\label{eq:delay_evidence}
	\end{equation}
	As given in Equation~\eqref{eq:delay_evidence}, $\mathcal{R}_{c_i}(t)$ is
	the set of roadside units currently assigned to controller $c_i$, and a
	threshold of $f{+}1$ independent roadside units is required before delay
	evidence contributes a penalty, preserving Byzantine fault tolerance
	against a single compromised unit falsely accusing a controller. The
	resulting controller trust score update is given in
	Equation~\eqref{eq:ctrl_trust_update}.
	\begin{equation}
		T_{c_i}^{(t+1)}{=}
		\begin{cases}
			\min(1,T_{c_i}^{(t)}{+}\Delta_r^{ctrl}) &
			\text{conflict evidence}{<}f{+}1
			\;\land\;
			\mathcal{E}_{delay}(c_i,t){<}f{+}1
			\\
			\max(0,T_{c_i}^{(t)}{-}\Delta_p^{ctrl}) &
			\text{conflict evidence}{\geq}f{+}1
			\;\lor\;
			\mathcal{E}_{delay}(c_i,t){\geq}f{+}1
		\end{cases}
		\label{eq:ctrl_trust_update}
	\end{equation}
	As given in Equation~\eqref{eq:ctrl_trust_update}, trust decrements when
	either conflict evidence or delay evidence reaches its threshold, so a
	timing manipulating controller is penalized through the delay channel
	even while every flow modification message it issues remains individually
	authorized. Distributed key generation among roadside unit ring members
	additionally removes the controller from the zero knowledge proof key
	plane entirely, so a compromised controller gains no cryptographic
	advantage from its position.
	
	
	\subsection{Dual Mode Detection Engine}\label{subsec:dual_mode}
	
	Detection is split across two cooperating stages so that no attack
	variant escapes coverage due to node capability. Figure~\ref{fig:dual_mode}
	shows the two parallel paths and their convergence at the blockchain
	layer.
	\begin{figure*}[!t]
		\centering
		\includegraphics[width=0.85\textwidth]{Figure_1_Final.eps}
		\caption{Dual Mode Framework Overview}
		\label{fig:dual_mode}
	\end{figure*}
	Resource constrained on board units run a lightweight rule engine
	evaluating only Signature S1 and the partial form of Signature S2 using
	mobility adjusted thresholds and hash based authentication, suspending
	forwarding on detection and escalating to the nearest roadside unit.
	Signature S1 must separate a genuine timing attack from ordinary
	congestion, which delays every traffic class rather than only safety
	critical packets. The condition formalising Signature S1 is given in
	Equation~\eqref{eq:sig_s1}.
	\begin{equation}
		\text{S1:}\quad
		\delta_p(v,r,t) > \bar{\delta}_r(t) + k\sigma_r(t)
		\;\land\;
		\text{Priority}(p) = \textsc{High}
		\;\land\;
		\delta_{best}(r,t) \leq \bar{\delta}_r(t) + k\sigma_r(t)
		\label{eq:sig_s1}
	\end{equation}
	As given in Equation~\eqref{eq:sig_s1}, the third conjunct requires best
	effort traffic at the same roadside unit to remain within baseline while
	high priority traffic does not, which is the selectivity property that an
	ordinary congestion induced delay, which elevates both classes together,
	cannot satisfy.
	
	Roadside units run the full pipeline, evaluating the remaining six
	signatures through blockchain observables and the cryptographic and
	witness layers of Sections~\ref{subsec:crypto_layer}
	and~\ref{subsec:witness}. Signatures S3 and S4, targeting slow flow
	table exhaustion, do not rely on installation rate at all, since a slow
	attacker deliberately operates within the legitimate rate distribution.
	Signature S3 instead fires on blockchain endorsement failure, given in
	Equation~\eqref{eq:rule_s3}.
	\begin{equation}
		f_{S3}(r,t) =
		\begin{cases}
			1 & f_{unauth}(r,t) = 1
			\;\land\; \nexists\,v :
			\text{flow}(r) \in \mathcal{F}_{active}(v) \\
			0 & \text{otherwise}
		\end{cases}
		\label{eq:rule_s3}
	\end{equation}
	As given in Equation~\eqref{eq:rule_s3}, $f_{unauth}{=}1$ is a binary,
	deterministic signal fixed by the endorsement mechanism of
	Section~\ref{subsec:blockchain}, achieving a zero false positive rate by
	construction rather than by threshold tuning. Signature S4 instead fires
	on table utilisation saturation, given in Equation~\eqref{eq:rule_s4}.
	\begin{equation}
		f_{S4}(r,t) =
		\begin{cases}
			1 & U_{TCAM}(r,t) > U_{thresh} \\
			0 & \text{otherwise}
		\end{cases}
		\;\text{with}\;
		v_{atk} = \arg\max_{v'}\,\lambda_{PI}(v',r,t)
		\label{eq:rule_s4}
	\end{equation}
	As given in Equation~\eqref{eq:rule_s4}, the attacking vehicle $v_{atk}$
	is attributed concurrently from the highest per source packet arrival
	rate but does not gate the detection decision itself, so an attribution
	failure never suppresses detection. Algorithm~\ref{alg:lrad} summarises
	the composite decision, which routes mitigation by attack plane: control
	plane detections update controller trust through
	Equation~\eqref{eq:ctrl_trust_update}, while data plane detections
	quarantine the source node directly.
	\begin{algorithm}[!t]
		\caption{Combined Lightweight and Full Detection}
		\label{alg:lrad}
		\begin{algorithmic}[1]
			\State \textbf{OBU:} evaluate S1, partial S2 on packet $p$; if either fires, suspend forwarding and escalate to the covering RSU
			\State \textbf{RSU:} evaluate full S2, S3 (Eq.~\ref{eq:rule_s3}), S4 (Eq.~\ref{eq:rule_s4}), S5 to S8 (Sections~\ref{subsec:crypto_layer}, \ref{subsec:witness})
			\State suppress LSTM contribution whenever S3 or S4 has fired this cycle
			\State $D_{RSU} \gets$ logical OR of all active RSU stage flags
			\If{$D_{RSU}$ and the firing signature is control plane}
			\State update controller trust, Eq.~\ref{eq:ctrl_trust_update}
			\ElsIf{$D_{RSU}$}
			\State quarantine the attributed source node
			\EndIf
		\end{algorithmic}
	\end{algorithm}
	
	
	\subsection{Post Quantum Cryptographic Layer}\label{subsec:crypto_layer}
	
	Every forwarded packet carries a per hop lattice based signature that
	aggregates across a multi hop path into a single verification rather
	than $n$ independent checks. Figure~\ref{fig:crypto_pipeline} shows the
	resulting pipeline together with the endorsement and commitment stages
	of Section~\ref{subsec:blockchain}.
	\begin{figure*}[!t]
		\centering
		\includegraphics[width=0.9\textwidth]{Figure_2_Final.eps}
		\caption{Cryptographic Packet Pipeline}
		\label{fig:crypto_pipeline}
	\end{figure*}
	The aggregate check is given in Equation~\eqref{eq:batch_verify}.
	\begin{equation}
		\textsc{BatchVerify}\!\left(
		\{\sigma_i\}, \{pk_i\}, \{m_i\}, \mathbf{r}
		\right) =
		\left[
		\sum_{i=1}^{n} r_i \mathbf{A}_i \mathbf{z}_i
		\stackrel{?}{=}
		\sum_{i=1}^{n} r_i
		\!\left(
		\mathbf{t}_i \mathbf{c}_i +
		\mathbf{w}'_i
		\right) \pmod{q}
		\right]
		\label{eq:batch_verify}
	\end{equation}
	As given in Equation~\eqref{eq:batch_verify}, a single Fiat Shamir
	challenged lattice equation accepts or rejects all $n$ per hop signatures
	at once, and any individual forgery perturbs the random linear
	combination with overwhelming probability, reducing the verification
	cost of a multi hop path from $n$ separate lattice checks to one
	combined operation.
	
	Timing compliance and next hop legitimacy are each certified by a
	transparent zero knowledge proof, a STARK, rather than a raw timestamp,
	so no verifier ever observes a vehicle's true delay or routing choice.
	Each forwarding node proves, without revealing its reception timestamp,
	forwarding timestamp, or its next hop's receipt signature, that the
	delay constraint and next hop authorisation were both satisfied.
	Verification of the two proofs is given in
	Equation~\eqref{eq:stark_verify}.
	\begin{equation}
		b_{delay,i} = \textsc{STARK.Verify}(\pi_{delay,i}, C_{delay}, H_{\text{SHA3-512}}) \in \{0,1\},
	\end{equation}
	\begin{equation}
		b_{hop,i} = \textsc{STARK.Verify}(\pi_{hop,i}, C_{hop}, H_{\text{SHA3-512}}) \in \{0,1\}
		\label{eq:stark_verify}
	\end{equation}
	As given in Equation~\eqref{eq:stark_verify}, $b_{delay,i}{=}0$ flags a
	node that delayed the packet beyond the safety bound, contributing to
	Signatures S1 and S2, while $b_{hop,i}{=}0$ is a corroborating
	cryptographic indicator for hidden forwarding rather than the primary
	one, since a compromised node controls its own proof generation and can
	produce a valid proof while covertly duplicating over the radio medium.
	These proofs alone cannot separate a genuine hidden forwarding duplicate
	from a bystander that merely overheard a broadcast transmission, since
	every neighbour within range receives every frame. Signature S5
	resolves this by requiring that the accusing node was the addressed
	recipient of the duplicate, given in Equation~\eqref{eq:sig_s5}.
	\begin{equation}
		\begin{aligned}
			\text{S5:}\quad
			& d' \notin \mathcal{P}(s,d) \\
			& \wedge\; \neg\,\textsc{BC.Endorsed}\big(\texttt{FlowMod}(s \to d')\big) \\
			& \wedge\; \textsc{ML-DSA-87.Verify}(
			\sigma_{copy}, pk_s, m_{copy}) = 0 \\
			& \wedge\; b_{hop}(u) = 0 \\
			& \wedge\; \textsc{Addressed}(p, d') = 1
		\end{aligned}
		\label{eq:sig_s5}
	\end{equation}
	As given in Equation~\eqref{eq:sig_s5}, the addressing conjunct is read
	directly from the link layer destination of the received frame rather
	than from any privileged knowledge of which nodes are compromised, and it
	is what allows the remaining conjuncts to distinguish a node that was
	sent an unauthorized duplicate from one that only overheard a legitimate
	transmission. Signatures S6 through S8 follow the same construction
	against their respective active and passive threat models.
	
	
	\subsection{Federated LSTM with Byzantine Robust Aggregation}\label{subsec:fed_lstm}
	
	Rule based and cryptographic signatures cover every attack with a
	deterministic observable, but passive hidden forwarding leaves no
	timing, volume, or cryptographic trace, since the retransmitted packet is
	a genuinely authentic copy. Detection of this residual class is
	delegated to a recurrent model trained across roadside units without
	centralizing raw traffic. Figure~\ref{fig:fed_lstm} shows the
	architecture together with the blockchain verified aggregation described
	below.
	\begin{figure*}[!t]
		\centering
		\includegraphics[width=0.9\textwidth]{Figure_3_Final.eps}
		\caption{Federated LSTM Architecture}
		\label{fig:fed_lstm}
	\end{figure*}
	Each roadside unit observes an eleven feature vector per cycle, given in
	Equation~\eqref{eq:lstm_input}.
	\begin{equation}
		\mathbf{x}_t^{(r)} = \left[
		\delta^{\max}_t,\;
		\lambda_{PI,t},\;
		U_{TCAM,t},\;
		\mathbb{1}[\pi_{delay}{=}\bot],\;
		\mathbb{1}[\pi_{hop}{=}\bot],\;
		\rho_t,\;
		\bar{v}_t,\;
		\log(1{+}N_{div,t}),\;
		\log(1{+}A^{\mathrm{RSU}}_{tp,t}),\;
		\log(1{+}R_{anom,t}),\;
		\mathbb{1}[\delta^{\max}_t{>}\Delta_{max}]
		\right]^T
		\label{eq:lstm_input}
	\end{equation}
	As given in Equation~\eqref{eq:lstm_input}, three features, the
	destination diversity count, the per roadside unit throughput asymmetry,
	and the reception log anomaly score, are transformed by $\log(1{+}x)$
	before entering the vector, since their benign resting value is small
	and constant while their value under attack spans several orders of
	magnitude; left untransformed, they would dominate every other feature
	after normalisation. The reception log anomaly score is the only feature
	that remains informative for the two passive variants, since every
	cryptographic and volumetric observable is unchanged when the original
	signed packet is merely retransmitted rather than modified.
	
	An autoencoder over this vector is trained first and its encoder is then
	frozen; a lightweight classification head is trained separately on the
	frozen encoder's final hidden state, given in Equation~\eqref{eq:lstm_cls}.
	\begin{equation}
		P_{atk}^{(k)}(t) =
		\sigma\!\left(\mathbf{w}_2^{\top}\,
		\mathrm{ReLU}\!\left(\mathbf{W}_1 \mathbf{h}^{(2)}_T + \mathbf{b}_1\right)
		+ b_2\right)
		\label{eq:lstm_cls}
	\end{equation}
	As given in Equation~\eqref{eq:lstm_cls}, the head is trained with binary
	cross entropy against injection side ground truth rather than against
	any quantity computed from the input vector itself, which keeps the
	label independent of the features the model consumes. Detection follows
	a per roadside unit threshold calibrated from that unit's own benign
	traffic under live operating conditions, given in
	Equation~\eqref{eq:lstm_cls_threshold}.
	\begin{equation}
		D_{LSTM}^{(k)}(t) = \mathbb{1}\!\left[P_{atk}^{(k)}(t) > \theta_{cls}^{(k)}\right],
		\qquad
		\theta_{cls}^{(k)} = \max\!\left(
		Q_{1-\alpha}\!\left(\mathcal{P}^{(k)}_{benign}\right),\; 0.05\right)
		\label{eq:lstm_cls_threshold}
	\end{equation}
	As given in Equation~\eqref{eq:lstm_cls_threshold}, the threshold is the
	empirical benign quantile of the roadside unit's own classifier output
	under live operating conditions rather than a value calibrated offline,
	since offline calibrated thresholds transfer poorly across the
	heterogeneous traffic each roadside unit observes once deployed.
	
	Local models are aggregated under the assumption that any roadside unit
	may be compromised. Algorithm~\ref{alg:brfa} summarises the four stage
	aggregation.
	\begin{algorithm}[!t]
		\caption{Byzantine Robust Federated Aggregation}
		\label{alg:brfa}
		\begin{algorithmic}[1]
			\State Trust gate: keep roadside units with trust at or above $T_{min}$; abort if fewer than $2f{+}1$ remain
			\State Hash verification: discard any submission whose weights do not match its committed hash
			\State Krum filter: compute the coordinate wise median $\tilde{\mathbf{W}}$ of the remaining submissions; discard any submission further than $\gamma$ from it
			\State Weighted aggregation: average the surviving submissions by sample count and commit the resulting hash to the blockchain
		\end{algorithmic}
	\end{algorithm}
	
	
	\subsection{Witness Based Passive Attack Detection}\label{subsec:witness}
	
	Passive hidden forwarding survives every mechanism described above,
	since the duplicated packet is a genuine, unmodified retransmission
	carrying a valid signature. Coverage of this residual case falls to
	nearby vehicles acting as witnesses, exploiting the fact that a
	broadcast transmission is overheard by every neighbour rather than only
	its intended recipient. A witness raises a duplication alert whenever it
	observes the same packet hash addressed to two distinct destinations,
	given in Equation~\eqref{eq:dup_alert_cond}.
	\begin{equation}
		\text{DA}_w(p){=}1\iff
		\exists\,dst{\neq}dst':
		H(p)\in\mathcal{L}_w(dst,W)
		\land H(p)\in\mathcal{L}_w(dst',W)
		\label{eq:dup_alert_cond}
	\end{equation}
	As given in Equation~\eqref{eq:dup_alert_cond}, this condition is
	impossible for legitimate, single path traffic, so any occurrence is
	itself evidence of duplication regardless of which node is ultimately
	found responsible.
	
	A single witness alert must never be sufficient to penalize a node,
	since one compromised or mistaken witness could otherwise accuse an
	innocent node at will. The smart contract penalty trigger therefore
	requires a distinct witness quorum, given in Equation~\eqref{eq:bft_penalty}.
	\begin{equation}
		\textsc{SC.PenalizeTrust}(v_i) \Leftarrow
		\Bigl|\bigl\{
		w : \exists\,(\alpha_w \text{ or } \beta_w)
		\text{ s.t. }
		\textsc{ML-DSA-87.Verify}(\cdot,pk_w,\cdot) = 1
		\;\land\; \text{event}(v_i,p,w)\text{ confirmed}
		\bigr\}\Bigr|
		\geq 2f{+}1
		\label{eq:bft_penalty}
	\end{equation}
	As given in Equation~\eqref{eq:bft_penalty}, the cardinality is taken
	over distinct witness vehicles that each independently submitted a
	verified alert for the same event, not over the number of alert
	messages, so a single vehicle resubmitting the same alert can never
	cross the threshold unilaterally.
	
	
	\subsection{Blockchain Architecture}\label{subsec:blockchain}
	
	All blockchain writes originate exclusively from roadside units, since
	permitting the controller itself to write would reduce the ledger's
	integrity guarantee to that of the very entity the architecture assumes
	may be compromised. Figure~\ref{fig:blockchain_arch} shows the resulting
	two tier structure.
	\begin{figure*}[!t]
		\centering
		\includegraphics[width=0.9\textwidth]{Figure_4_Final.eps}
		\caption{Blockchain Architecture and Trust Management}
		\label{fig:blockchain_arch}
	\end{figure*}
	A local chain among neighbouring roadside units commits flow
	modification endorsements and packet receipt logs under a low latency
	consensus tolerant of up to a third dishonest participants, while a
	global chain periodically anchors each local chain's state for
	auditability. A flow modification is only ever installed once it carries
	$f{+}1$ independent roadside unit endorsements, which is what allows
	Signature S3 of Equation~\eqref{eq:rule_s3} to detect an unauthorized
	installation without ever consulting the controller's own account of its
	actions.
	
	
	\section{Evaluation Setup}\label{evaluation_setup}
	
	All three experiments in Section~\ref{results} use a Manhattan, New York
	City road network extracted from OpenStreetMap rather than the Los
	Angeles trace used in the companion single class evaluations, at the
	same vehicle counts and speed settings. Manhattan's higher intersection
	density produces more frequent, more aggressive flow modification churn,
	giving a more demanding camouflage environment for the slow flow table
	exhaustion variants than a lower density suburban grid provides.
	Table~\ref{tab:sim_settings} lists the settings required for
	replication; tuning constants that are not central to the proposed
	design, together with their full sensitivity sweeps, are reported in the
	companion papers.
	
	NEXUS is compared against three baselines: TAP for timing based
	selective delay, SFTO Guard for slow flow table exhaustion, and FADE for
	forwarding anomaly detection, none of which were designed to operate
	against more than one attack class at a time. Every baseline is
	evaluated under the identical joint, eight variant traffic used to
	evaluate NEXUS itself, which is what exposes the joint detection gap
	reported in Section~\ref{results}.
	
	\begin{table}[!t]
		\centering
		\caption{Simulation Settings}
		\label{tab:sim_settings}
		\renewcommand{\arraystretch}{1.2}
		\footnotesize
		\begin{tabular}{|>{\raggedright\arraybackslash}p{3.2cm}|>{\raggedright\arraybackslash}p{4.6cm}|}
			\hline
			\textbf{Parameter} & \textbf{Value} \\
			\hline
			Mobility trace & Manhattan, New York City (OpenStreetMap) \\
			\hline
			Network scale (default) & 200 vehicles, 64 RSUs (8$\times$8 grid), 4 controllers \\
			\hline
			Data / control plane & IEEE 802.11p DSRC / Ethernet CSMA \\
			\hline
			Attack configuration & 8 variants active simultaneously, penetration swept per experiment, 5 seeds \\
			\hline
			Detection architecture & Dual mode rule engine (OBU/RSU) plus federated classifier head, gated off for S3/S4 \\
			\hline
			Mobility baseline, EWMA & $\beta = 0.95$ ($N_{eff}{=}20$ cycles); $k = 3$ \\
			\hline
			LSTM classifier threshold & Per RSU empirical benign quantile, $\alpha = 0.01$, floor $0.05$ \\
			\hline
			Byzantine aggregation & Krum threshold $\gamma = 2.0$; trust gate $T_{min} = 0.50$ \\
			\hline
			Cryptographic layer & ML-DSA-87 signatures; batch size $B = 15$ packets \\
			\hline
			Blockchain endorsement & $f{+}1$ RSU endorsements per FlowMod; $f = 1$ \\
			\hline
			Witness quorum & $2f{+}1{=}3$ distinct witnesses; window $W{=}10$s \\
			\hline
			Duplication / volume windows & S6 duplication window 30s; S7 tumbling volume window 10s \\
			\hline
			Baselines & TAP, SFTO Guard, FADE \\
			\hline
			Ground truth, attribution & Latched for hidden forwarding variants; suspect attributed \\
			\hline
		\end{tabular}
	\end{table}
	
	Detection quality throughout Section~\ref{results} is reported as macro
	averaged MCC across the eight attack variants, with the pooled figure
	shown alongside for transparency only, following the convention
	established in the companion papers. Two further metrics, defined in
	full in the companion papers, are reported alongside it: threshold
	violation rate, the share of safety critical packets that exceed the
	selective time delay threshold of Section~\ref{subsec:dual_mode}, and
	unauthorized copy rate, the share of hidden forwarding duplicates that
	reach an unauthorized destination undetected. Mitigation latency is
	reported as the pair described in Section~\ref{discussion}, a
	prevention rate and a mean latency restricted to attackers that acted
	before containment, rather than as a single mean, and end to end
	latency is the delivery time from packet generation to final delivery
	inclusive of every detection and consensus step on its path.
	
	
	\section{Results}\label{results}
	
	% All quantitative values in this section are placeholders pending the
	% Manhattan trace runs under the finalised configuration of Table
	% \ref{tab:sim_settings}. Replace every \TBD with the measured figure
	% once available; the analysis prose states the mechanism each trend is
	% expected to follow and should be checked against, not discarded,
	% once real numbers land.
	%
	% N3 specifically cannot be executed as scoped until three prerequisite
	% fixes land: (i) the No Zero Trust arm must be made compromisable below
	% p=100 percent, since a single controller currently cannot be attacked
	% at all under the existing allocation logic; (ii) a clean toggle for
	% zero trust off is needed that does not also disable signing, since the
	% only existing switch conflates the two; (v) a counter is needed at the
	% enforcement call sites to confirm a quarantine actually blocked an
	% action rather than only being recorded. None of this affects the
	% experimental description below, only whether it can be run yet.
	
	\subsection{Joint Concurrent Multi Class Detection Under Varying Attack Penetration}\label{subsec:n1}
	
	All eight attack variants are active simultaneously at every penetration
	level $p \in \{0,20,40,60,80,100\}\%$, which is the central
	differentiator from the companion papers, both of which evaluate a
	single attack class at a time. Figure~\ref{fig:n1_mcc} reports macro
	MCC against penetration for NEXUS and the three baselines.
	\begin{figure}[!t]
		\centering
		\fbox{\begin{minipage}[c][4.2cm]{0.92\linewidth}
				\centering
				Figure placeholder: macro MCC versus attack penetration, four lines
				(NEXUS, TAP, SFTO Guard, FADE), $p \in \{0,20,40,60,80,100\}\%$.
		\end{minipage}}
		\caption{Detection Quality Under Joint Attack Penetration}
		\label{fig:n1_mcc}
	\end{figure}
	Every baseline is structurally blind to attack classes outside its own
	design scope: TAP produces no signal against any hidden forwarding
	variant, SFTO Guard produces no signal outside slow flow table
	exhaustion, and FADE produces no signal against selective time delay.
	Under a joint, eight variant attack, this blindness is not merely a
	missed detection on the unseen classes, it also degrades detection on
	the classes each baseline was designed for, since FADE's flow
	conservation baseline is itself perturbed by concurrent selective delay,
	and TAP's timing baseline is perturbed by concurrent flow table
	contention. Macro MCC for every baseline is therefore expected to fall
	well below its single class figure reported in the companion papers,
	while NEXUS, whose eight signatures and federated classifier run
	concurrently by design, is expected to hold macro MCC closer to its own
	single class figures as penetration rises.
	
	Figure~\ref{fig:n1_tvr_ucr} reports the threshold violation rate and
	unauthorized copy rate on the same penetration axis.
	\begin{figure}[!t]
		\centering
		\fbox{\begin{minipage}[c][4.2cm]{0.92\linewidth}
				\centering
				Figure placeholder: TVR and UCR versus attack penetration, dual
				axis, NEXUS and the three baselines.
		\end{minipage}}
		\caption{Threshold Violation and Unauthorized Copy Under Joint Attack}
		\label{fig:n1_tvr_ucr}
	\end{figure}
	Unauthorized copy rate is expected to separate NEXUS most clearly from
	FADE specifically, since a hidden forwarding baseline evaluated while
	timing attacks are simultaneously active must still classify duplication
	correctly under a degraded timing baseline, whereas NEXUS's hidden
	forwarding signatures and witness layer of
	Section~\ref{subsec:witness} do not depend on any timing observable at
	all. At $p{=}\TBD\%$, macro MCC for NEXUS is $\TBD$ against $\TBD$ for
	the strongest baseline; unauthorized copy rate for NEXUS remains at
	$\TBD\%$ against $\TBD\%$ for FADE.
	
	
	\subsection{Scalability and Cryptographic Overhead Under Concurrent Attacks}\label{subsec:n2}
	
	Network scale is varied as $N_v \in \{100,200,300,400\}$ vehicles at a
	fixed $40\%$ joint penetration, holding the roadside unit grid and
	controller count constant. Figure~\ref{fig:n2_scale} reports end to end
	latency against the three baselines together with the two cryptographic
	overhead terms that only NEXUS incurs.
	\begin{figure}[!t]
		\centering
		\fbox{\begin{minipage}[c][4.2cm]{0.92\linewidth}
				\centering
				Figure placeholder, two panels: (a) end to end latency versus
				vehicle count, NEXUS and the three baselines; (b) batch
				verification time and endorsement consensus latency versus
				vehicle count, NEXUS only.
		\end{minipage}}
		\caption{Scalability and Cryptographic Overhead}
		\label{fig:n2_scale}
	\end{figure}
	Endorsement consensus latency is expected to grow sub linearly with
	vehicle count, since the endorsement quorum of Section~\ref{subsec:blockchain}
	is sized by the fixed roadside unit count rather than by the vehicle
	population, so only per packet batch verification cost, amortized
	across a batch by Equation~\eqref{eq:batch_verify}, tracks vehicle
	driven traffic growth directly. Baselines are expected to report lower
	end to end latency at every vehicle count, but this reflects the absence
	of any cryptographic or blockchain overhead in their design rather than
	a genuine efficiency advantage under an equal integrity requirement. At
	$N_v{=}400$, end to end latency for NEXUS is $\TBD$ms against a
	$100$ms safety bound, and endorsement consensus latency is $\TBD$ms
	against $\TBD$ms at $N_v{=}100$.
	
	
	\subsection{System Level Layer Ablation}\label{subsec:n3}
	
	Four configurations isolate whole architectural layers rather than
	individual components, in contrast to the component level ablations
	reported in the companion papers: Full NEXUS; No Crypto, removing
	ML-DSA-87 and STARK attestation entirely; No LSTM, removing the
	federated classifier and its aggregation; and No Zero Trust, replacing
	the multi controller architecture of Section~\ref{subsec:ctrl_trust}
	with a single, permanently trusted controller. All four run under the
	joint, eight variant attack at $40\%$ penetration.
	\begin{figure}[!t]
		\centering
		\fbox{\begin{minipage}[c][4.2cm]{0.92\linewidth}
				\centering
				Figure placeholder: radar chart, four spokes (macro MCC,
				unauthorized copy rate, mitigation prevention rate, end to end
				latency), one polygon per configuration.
		\end{minipage}}
		\caption{System Level Layer Ablation}
		\label{fig:n3_radar}
	\end{figure}
	Removing the cryptographic layer is expected to collapse containment of
	the active hidden forwarding variants specifically, since signature
	verification failure is the primary evidence Signature S5 and its
	active counterpart test for; it additionally removes the signing
	mechanism the witness layer uses to attribute an alert non repudiably,
	so passive hidden forwarding coverage may degrade as well rather than
	surviving on the classifier alone. Removing the federated classifier is
	expected to reduce, but not eliminate, coverage of the two passive
	variants, since the witness quorum of Equation~\eqref{eq:bft_penalty}
	operates independently of the classifier and continues to provide a
	detection path once it is removed. Removing multi controller zero trust
	leaves detection quality nominally unchanged below the compromise
	threshold, but failover latency becomes undefined, since a single
	controller design has no second controller to fail over to, and macro
	MCC is expected to collapse once that one controller is compromised,
	with no independent recovery path available.
	
	Table~\ref{tab:n3_variant} reports macro MCC for each attack variant
	under each configuration, together with unauthorized copy rate for the
	full system.
	\begin{table}[!t]
		\centering
		\caption{Per Variant MCC Under System Level Ablation}
		\label{tab:n3_variant}
		\renewcommand{\arraystretch}{1.15}
		\footnotesize
		\begin{tabular}{|c|c|c|c|c|c|}
			\hline
			\textbf{Variant} & \textbf{Full} & \textbf{No Crypto} & \textbf{No LSTM} & \textbf{No Zero Trust} & \textbf{UCR (Full)} \\
			\hline
			S1 & \TBD & \TBD & \TBD & \TBD & --- \\
			\hline
			S2 & \TBD & \TBD & \TBD & \TBD & --- \\
			\hline
			S3 & \TBD & \TBD & \TBD & \TBD & --- \\
			\hline
			S4 & \TBD & \TBD & \TBD & \TBD & --- \\
			\hline
			S5 & \TBD & \TBD & \TBD & \TBD & \TBD \\
			\hline
			S6 & \TBD & \TBD & \TBD & \TBD & \TBD \\
			\hline
			S7 & \TBD & \TBD & \TBD & \TBD & \TBD \\
			\hline
			S8 & \TBD & \TBD & \TBD & \TBD & \TBD \\
			\hline
		\end{tabular}
	\end{table}
	Under controller compromise at $p \geq 33\%$, macro MCC for
	No Zero Trust falls to $\TBD$ against $\TBD$ for Full NEXUS under the
	identical compromise scenario, and mitigation prevention rate falls to
	$\TBD\%$.
	
	
	
	\section{Discussion}\label{discussion}
	
	Mitigation latency is reported throughout this paper as a pair rather
	than a single mean: a prevention rate, the share of quarantined
	attackers stopped before their first attack action, and a mean latency
	computed only over attackers that acted before containment. A single
	mean over both populations moves in the wrong direction as containment
	improves, since an attacker stopped before acting contributes no
	latency at all and would otherwise be discarded from, rather than
	credited to, the average.
	
	\subsection{Statistical Significance}
	Every experiment in Section~\ref{results} is run across five seeds.
	The comparison between NEXUS and the strongest baseline at each
	penetration level in Experiment N1 is tested with a paired test across
	seeds rather than a comparison of seed averages, since the same five
	mobility realisations are shared across every method. At
	$p{=}\TBD\%$, the macro MCC difference is $\TBD$ with $\TBD$
	($n{=}5$ paired seeds); this is reported once the Manhattan trace runs
	are complete and is not to be read as significant or otherwise until
	then.
	
	\subsection{Strengths and Weaknesses}
	The central strength demonstrated in Section~\ref{results} is joint
	coverage: every signature, the cryptographic layer, and the federated
	classifier operate concurrently by construction, so detection of one
	attack class is not disrupted by the simultaneous presence of the
	other, which Experiment N1 is designed specifically to expose as a
	failure mode in every baseline. The delay evidence channel of
	Equation~\eqref{eq:ctrl_trust_update} closes a documented gap in
	endorsement only controller trust, and post quantum attestation is
	achieved with batch verification that keeps per packet overhead sub
	linear in path length.
	
	Three weaknesses are structural rather than incidental. First, the
	logical OR across eight signatures and a classifier that composes
	overall detection trades false positive rate for coverage, so the
	composite decision is necessarily noisier than any single constituent
	detector taken alone. Second, saturation based signatures such as S4
	carry an inherent temporal lag on slow building attacks, since ground
	truth is established at the first attack conforming action while the
	signature itself requires accumulated saturation, a property of the
	detection mechanism rather than a failure of it. Third, passive hidden
	forwarding coverage depends on witness density, so the $2f{+}1$ quorum
	of Equation~\eqref{eq:bft_penalty} becomes harder to satisfy exactly
	where vehicle density is lowest.
	
	\subsection{Runtime and Space Complexity}
	Aggregate signature verification is linear in path length and Byzantine
	robust aggregation is quadratic in the number of participating roadside
	units, with the roadside unit term dominant at the network scale
	evaluated in Section~\ref{results}. Appendix~\ref{appendix} states and
	proves the formal bounds.
	
	\subsection{Scalability}
	Consistent with Experiment N2, end to end latency is expected to remain
	within the safety bound as vehicle count grows, since the dominant
	consensus cost is fixed by roadside unit count rather than by vehicle
	population. The complexity result of Appendix~\ref{appendix} indicates
	that the binding cost as the network grows further is the quadratic
	aggregation term, which tracks roadside unit count rather than vehicle
	count, and is therefore not exercised by Experiment N2's scaling axis.
	
	\subsection{Deployment Feasibility}
	The dual mode design of Section~\ref{subsec:dual_mode} places
	post quantum signing, zero knowledge proof generation, and classifier
	inference exclusively at roadside unit edge nodes, while on board units
	run only hash based authentication, so the heavier computational
	burden falls on infrastructure hardware rather than on the vehicle
	itself. Distributed key generation is a one time ceremony at network
	setup, and key rotation is triggered only on roadside unit revocation,
	so neither recurs on the critical path of ordinary operation.
	
	\subsection{Convergence}
	Byzantine robust aggregation converges toward the honest aggregate
	under the bounded malicious fraction assumed throughout this paper.
	Appendix~\ref{appendix} states and proves the bound; empirical support
	under the shared aggregation ablation is reported in the companion
	papers.
	
	\subsection{Security Analysis}
	False accusation of a benign forwarding node under Signature S5 and its
	counterparts requires either breaking ML-DSA-87 unforgeability or
	compromising the physical layer addressing the addressed conjunct of
	Equation~\eqref{eq:sig_s5} is read from, and the controller trust
	channel of Equation~\eqref{eq:ctrl_trust_update} cannot be biased by a
	single compromised roadside unit. Appendix~\ref{appendix} states and
	proves both claims, together with the passive attack limitation that
	motivates the witness layer of Section~\ref{subsec:witness}.
	
	
	\section{Conclusion}\label{conclusion}
	
	Selective time delay and hidden forwarding attacks in software defined
	vehicular networks have been addressed by separate defenses, each built
	around a threat model that assumes a trustworthy controller and a
	single attack class. NEXUS removes both assumptions at once: a flat,
	registered controller set with no permanently trusted member, a dual
	mode detection engine spanning eight signatures, a post quantum
	cryptographic layer with a controller inclusive trust channel, a
	federated classifier trained without centralising raw traffic, and a
	witness layer covering what cryptographic and volumetric evidence
	structurally cannot. Under joint, simultaneous attack across all eight
	variants, NEXUS is shown in Section~\ref{results} to sustain macro MCC
	of $\TBD$ against $\TBD$ for the strongest single class baseline at
	the same penetration, while every baseline exhibits the coverage
	collapse this paper set out to demonstrate. The main limitation is
	architectural rather than incidental: composing eight signatures and a
	classifier through a logical OR trades false positive rate for the
	coverage that joint detection requires, and this trade is not removed
	by any of the individual fixes reported here.
	
	Two companion papers extend the single class evaluation this paper
	condenses: PHANTOM develops the selective time delay signatures and
	mobility adjusted thresholds in full formal detail, and VANGUARD-HF
	does the same for hidden forwarding, the federated classifier, and the
	witness mechanism. Future work includes replacing the covering roadside
	unit proxy used to contain compromised vehicles with a per vehicle
	trust score, re running the Byzantine robust aggregation ablation on
	the corrected eleven feature input once it is available across both
	companion papers, and a density adaptive witness quorum that lowers
	$f$ where vehicle density is too sparse to sustain $2f{+}1$ distinct
	witnesses, directly addressing the limitation identified in
	Section~\ref{discussion}.
	
	\subsection{Acknowledgments}
	None
	
	\subsection{Conflict of interest}
	\noindent Authors declare no conflict of interest.
	
	\section[\appendixname~\thesection]{Appendix}\label{appendix}
	
	\subsection{Complexity Analysis}\label{app:complexity}
	
	Table~\ref{tab:complexity} states the per operation cost of the four
	dominant mechanisms in the pipeline.
	\begin{table}[!t]
		\centering
		\caption{Runtime and Space Complexity}
		\label{tab:complexity}
		\renewcommand{\arraystretch}{1.2}
		\footnotesize
		\begin{tabular}{|l|c|c|}
			\hline
			\textbf{Operation} & \textbf{Time} & \textbf{Space} \\
			\hline
			BatchVerify, Eq.~\eqref{eq:batch_verify} & $O(n)$ & $O(n)$ \\
			\hline
			STARK.Verify per proof & $O(\log n)$ & $O(1)$ \\
			\hline
			Krum filter, Algorithm~\ref{alg:brfa} & $O(K^2)$ & $O(Kd)$ \\
			\hline
			Endorsement, Section~\ref{subsec:blockchain} & $O(f{+}1)$ & $O(f{+}1)$ \\
			\hline
		\end{tabular}
	\end{table}
	As given in Table~\ref{tab:complexity}, $n$ is the path length in hops,
	$K$ is the number of participating roadside units, $d$ is the model
	weight dimension, and $f$ is the Byzantine fault parameter.
	
	\begin{lemma}
		\label{lem:complexity}
		The per packet verification cost of the cryptographic layer is
		$O(n)$ in path length $n$, and the per round cost of Byzantine
		robust aggregation is $O(K^2)$ in the number of participating
		roadside units $K$, independently of $n$.
	\end{lemma}
	\begin{proof}
		BatchVerify of Equation~\eqref{eq:batch_verify} replaces $n$
		independent lattice checks with one combined matrix accumulation
		over the $n$ per hop signatures; the accumulation itself is
		$O(n)$, so the asymptotic cost is unchanged while the constant
		factor is reduced relative to $n$ separate verifications.
		STARK.Verify is polylogarithmic in the size of the statement being
		proven for a transparent, FRI based proof system, so it does not
		grow with $n$ beyond the single proof already accounted for per
		hop. The Krum filter of Algorithm~\ref{alg:brfa} computes a
		pairwise distance between every pair of the $K$ submitted models
		to identify the coordinate wise median neighbourhood, which is
		$O(K^2)$ regardless of $n$.
		\qed
	\end{proof}
	
	\begin{corollary}
		At the default network scale of Table~\ref{tab:sim_settings}, with
		$K{=}64$ roadside units, the Krum filter dominates the total per
		round cost whenever the evaluated path length satisfies
		$n \ll K^2$, which holds for every path length reachable within
		the eight by eight grid deployment.
	\end{corollary}
	
	\begin{remark}[Limitation]
		The quadratic term in Lemma~\ref{lem:complexity} is not exercised
		by Experiment N2, which scales vehicle count rather than roadside
		unit count; a deployment that instead scales the roadside unit
		grid would need the aggregation round cost re evaluated at the
		larger $K$.
	\end{remark}
	
	
	\subsection{Convergence Analysis}\label{app:convergence}
	
	\begin{lemma}
		\label{lem:convergence}
		Let $\rho_{mal}$ be the fraction of roadside units submitting
		adversarial weight updates in a given aggregation round. If
		$\rho_{mal} < 1/3$, the output of Algorithm~\ref{alg:brfa} lies
		within a bounded distance of the aggregate that would be produced
		by the honest submissions alone, for a Krum threshold $\gamma$
		chosen at least as large as the spread of the honest population.
	\end{lemma}
	\begin{proof}
		The trust gate and hash verification steps of
		Algorithm~\ref{alg:brfa} admit only submissions from roadside
		units above the operational trust threshold with a weight matching
		their committed hash, so any adversarial submission reaching the
		Krum filter carries no advantage from forged identity or a
		mismatched payload. The coordinate wise median
		$\tilde{\mathbf{W}}$ used by the Krum filter is itself robust to
		up to, but not including, half of the remaining population being
		adversarial, so for $\rho_{mal}<1/3$ the median stays within the
		convex spread of the honest submissions. Any submission further
		than $\gamma$ from $\tilde{\mathbf{W}}$ is excluded before
		averaging; once $\gamma$ is set at least as large as the honest
		population's own spread around the median, every honest submission
		is retained while an adversarial submission distant enough to bias
		the average by more than the stated bound is excluded, and the
		weighted aggregation step of Algorithm~\ref{alg:brfa} therefore stays
		within a bounded distance of the honest only aggregate.
		\qed
	\end{proof}
	
	\begin{remark}[Limitation]
		Lemma~\ref{lem:convergence} bounds the deviation from the honest
		aggregate per round rather than guaranteeing convergence of the
		underlying model to a global optimum, which additionally depends
		on the non convexity of the encoder trained in
		Section~\ref{subsec:fed_lstm}; empirical support for the stated
		bound under the shared Byzantine robust aggregation ablation is
		reported in the companion papers rather than repeated here.
	\end{remark}
	
	
	\subsection{Security Analysis}\label{app:security}
	
	\begin{lemma}
		\label{lem:delay_evidence_bft}
		No single compromised roadside unit can cause a controller to be
		penalised through the delay evidence channel of
		Equation~\eqref{eq:delay_evidence}.
	\end{lemma}
	\begin{proof}
		Equation~\eqref{eq:ctrl_trust_update} decrements controller trust
		on the delay evidence channel only once
		$\mathcal{E}_{delay}(c_i,t) \geq f{+}1$, and
		Equation~\eqref{eq:delay_evidence} defines this quantity as a
		cardinality over distinct roadside units independently reporting
		Signature S1 against the same controller. A single compromised
		roadside unit can contribute at most one to this count, which is
		insufficient to reach the $f{+}1$ threshold whenever $f \geq 1$,
		so a lone false report cannot by itself trigger a trust penalty.
		\qed
	\end{proof}
	
	\begin{lemma}
		\label{lem:s5_unforgeability}
		Falsely accusing a benign forwarding node under Signature S5 of
		Equation~\eqref{eq:sig_s5} requires either breaking ML-DSA-87
		unforgeability or compromising the link layer addressing that the
		addressed conjunct is read from.
	\end{lemma}
	\begin{proof}
		A benign forwarding node relays a packet under a signature chain
		originating from its genuine sender, so
		$\textsc{ML-DSA-87.Verify}(\sigma_{copy},pk_s,m_{copy})=1$ for any
		copy the accusing node genuinely received; producing a copy that
		fails this check while attributing it to a benign node's action
		requires forging a valid looking signature under the sender's
		public key, which is precluded by the module lattice hardness
		assumption underlying ML-DSA-87 unforgeability. The addressing
		conjunct $\textsc{Addressed}(p,d')=1$ is read directly from the
		link layer destination of the frame the accusing node itself
		received, independent of the signed next hop, so satisfying it
		without genuinely having been addressed the duplicate requires
		control over the physical channel rather than over any
		cryptographic or protocol state.
		\qed
	\end{proof}
	
	\begin{corollary}
		Since Equation~\eqref{eq:bft_penalty} additionally requires
		$2f{+}1$ distinct witnesses to independently confirm an event
		before a trust penalty is applied, an adversary must both defeat
		Lemma~\ref{lem:s5_unforgeability} and collude across $2f{+}1$
		witnesses to falsely penalise a node through the witness path,
		which exceeds the Byzantine bound assumed throughout this paper.
	\end{corollary}
	
	\begin{remark}[Limitation]
		Lemma~\ref{lem:s5_unforgeability} does not extend to the two
		passive variants, since a passively retransmitted packet is a
		genuinely authentic copy that satisfies every cryptographic check
		by construction. This is precisely why passive hidden forwarding
		is covered by the independent witness layer of
		Section~\ref{subsec:witness} rather than by Signature S5 or S6
		alone, and it is the one class of accusation this appendix does
		not reduce to a cryptographic hardness assumption.
	\end{remark}
	
	
	
	\begin{thebibliography}{9}
		
		% NOTE: bibliographic details below were checked against publicly
		% indexed records where marked [verified]. Entries marked
		% [confirm before submission] rest on recollection of the source
		% material used earlier in this project and must be checked against
		% the publisher record before this paper is submitted, consistent
		% with the reference integrity standard already applied to the
		% companion papers.
		
		\bibitem{Pascoal2017} % [verified]
		T.~A. Pascoal, Y.~G. Dantas, I.~E. Fonseca, and V.~Nigam, ``Slow TCAM exhaustion DDoS attack,'' in \emph{Proc. 32nd IFIP Int. Conf. ICT Syst. Secur. Privacy Prot. (IFIP SEC)}, Rome, Italy, 2017, pp. 17--31.
		
		\bibitem{Arsalan2018} % [verified]
		A. Arsalan and R.~A. Rehman, ``Prevention of timing attack in software defined named data network with VANETs,'' in \emph{Proc. Int. Conf. Front. Inf. Technol. (FIT)}, Islamabad, Pakistan, 2018, pp. 247--252.
		
		\bibitem{Mohammadi2016} % [verified]
		A.~A. Mohammadi, P. Kazemian, and M. Pakravan, ``Detecting malicious packet drops and misroutings using header space analysis,'' in \emph{Proc. 8th Int. Symp. Telecommun. (IST)}, Tehran, Iran, 2016, pp. 521--526.
		
		\bibitem{Zhang2021FADE} % [confirm before submission]
		L. Zhang \emph{et al.}, ``FADE: Detecting forwarding anomalies in software defined networks via flow conservation,'' in \emph{Proc. IEEE Conf. Comput. Commun. (INFOCOM) Workshops}, 2021.
		
		\bibitem{Nayak2022} % [verified]
		R.~P. Nayak, S. Sethi, S.~K. Bhoi, D. Mohapatra, R.~R. Sahoo, P.~K. Sharma, and D. Puthal, ``TFMD-SDVN: A trust framework for misbehavior detection in the edge of software defined vehicular network,'' \emph{J. Supercomput.}, vol.~78, no.~6, pp. 7948--7981, 2022.
		
		\bibitem{Liu2021BFLVEC} % [venue confirmed, page range to confirm]
		H. Liu \emph{et al.}, ``Blockchain and federated learning for collaborative intrusion detection in vehicular edge computing,'' \emph{IEEE Trans. Veh. Technol.}, vol.~70, no.~6, pp. 6073--6084, 2021.
		
		\bibitem{Krishnan2023} % [confirm before submission]
		S. Krishnan \emph{et al.}, ``FTISCON: Preserving flow table integrity in software defined networks through smart contracts,'' in \emph{Proc. IEEE Int. Conf. Blockchain}, 2023.
		
		\bibitem{ElDalahmeh2025} % [confirm before submission]
		M. El-Dalahmeh and M. Alia, ``A lattice based post-quantum vehicle to roadside authentication protocol for connected vehicles,'' \emph{IEEE Access}, 2025.
		
		% TODO: add companion paper citations once PHANTOM and VANGUARD-HF
		% have assigned manuscript identifiers, cited here as under review.
		
	\end{thebibliography}
	
	\begin{IEEEbiography}[{\includegraphics[width=1in,height=1.25in,clip,keepaspectratio]{pictures/Nilmantha}}]{Patikiri Arachchige Don Shehan Nilmantha Wijesekara} received his B.Sc. Engineering degree (First Class Honours) specialising in Electrical and Information Engineering from the University of Ruhuna, Sri Lanka, in 2017, where he was awarded six academic distinctions, including two gold medals and a scholarship. He subsequently completed his Ph.D. in Computer Engineering at the same university, earning a further three gold medals for his doctoral work on networking, machine learning, and blockchain. He has authored more than 50 journal papers and holds an h-index of 15. Alongside his core work, he maintains research interests in biomedical engineering and the social sciences, particularly nutrition and mental health.
	\end{IEEEbiography}
	% TODO: add a second IEEEbiography block per additional author once the
	% author list is confirmed.
	
\end{document}